% Chapter 5: Theoretical Foundation
\section{Fundamentação Teórica}

Esta seção apresenta a base teórica que sustenta o projeto, abrangendo os principais conceitos relacionados a jogos sérios, gamificação e teorias de aprendizagem aplicadas ao design de jogos.

\subsection{Jogos Sérios e Gamificação}
Jogos sérios são jogos projetados para um propósito primário que não é o puro entretenimento. Conforme discutido por autores como Miyamoto \cite{conference_example}, o design de jogos pode ser aplicado a diversos contextos, incluindo treinamento, simulação e educação. A gamificação, por sua vez, refere-se ao uso de elementos e mecânicas de jogos em contextos de não-jogo para aumentar o engajamento e a motivação dos usuários.

\subsection{Teorias da Aprendizagem}
O design de "Alchemon" será influenciado por teorias da aprendizagem como o Construtivismo e o Conectivismo. O Construtivismo, popularizado por Piaget, postula que o aprendizado é um processo ativo de construção de novos conhecimentos com base em experiências anteriores. Um jogo que permite a exploração e a experimentação, como o proposto, alinha-se perfeitamente com esta visão.

O design de objetos e interfaces também deve considerar os princípios de usabilidade e design centrado no usuário, para garantir que a ferramenta seja intuitiva e fácil de usar, como Norman descreve em seu livro seminal \cite{book_example}. A clareza da interface é crucial para não criar barreiras ao aprendizado.

\begin{figure}[h]
    \centering
    % \includegraphics[width=0.9\columnwidth]{placeholder2}
    \caption{Exemplo de um ciclo de feedback em um jogo sério, onde o jogador age, recebe feedback e aprende.}
    \label{fig:feedback_loop}
\end{figure}
